\documentclass[Chp3.tex]{subfiles}

\begin{document}


\section{Conclusion}
\label{sec:chapter3-conclusion}

\par We have reached an alternative measure of wealth inequality, $I^{A}$. The measure both utilizes the notion of representative wealth levels and preserves the mean-independence property that many statistical measures of inequality have. The goal of this exercise is to apply it in the context of inequality in racial wealth. There, the story is generally told using the aforementioned statistical measures of inequality (most notable, the \textit{black-white median wealth gap}).

\par With the arguments made in section $3$, it seems reasonable to suspect that a state space such as $S=\{\text{black, white}\}$ should largely inform our measure of wealth inequality via some ranking over wealth distributions. Historical episodes such as enslavement and the era of Jim Crows laws suggest that race may be a socially-relevant feature in characterizing the level of inequality in a given wealth distribution. The main issue is whether or not the effects of these historical episodes persist over time, as well as how explanatory they are for racial differences in wealth.

\par The computational analysis in Subsection~\ref{subsec:compute-measures} delivers a clear result: the objective benchmark understates wealth inequality relative to the ambiguity-based measure. In SCF data (1989--2022), this comparison is robust for both net wealth and gross assets, and the resulting inequality series follow the same broad time pattern as the Black-White median wealth-gap diagnostics; the restricted sample is higher and smoother, while the shifted sample is more time-varying.

\end{document}
